% TOP_PROBLEM: {"id": 3162, "title": "The circular embedding conjecture", "category_id": 3, "category": {"id": 3, "name": "graph_theory", "display_name": "Graph Theory", "description": "Problems involving graphs, networks, and their properties.", "slug": "graph-theory", "order_index": 3, "created_at": "2026-07-31T15:26:25.671Z"}, "status": "open", "problem_number": "OPG-138", "set_id": 12}
% FIRST_POSTED: 2026-10-01
% ATTEMPT_STATUS: unresolved
% UnsolvedMath ID 3162; source catalogue code OPG-138
% !TeX program = lualatex
% GPT-6 Astra Ultra research attempt; independent partial work, 2026-10-01.
\documentclass[11pt,a4paper]{article}
\usepackage[margin=25mm]{geometry}
\usepackage{fontspec}
\setmainfont{lmroman10-regular.otf}[BoldFont=lmroman10-bold.otf,
 ItalicFont=lmroman10-italic.otf,BoldItalicFont=lmroman10-bolditalic.otf]
\usepackage{amsmath,amssymb,mathtools}
\usepackage{unicode-math}
\setmathfont{latinmodern-math.otf}
\AtBeginDocument{\let\setminus\smallsetminus}
\usepackage{xurl,hyperref}
\hypersetup{colorlinks=true,urlcolor=blue}
\setcounter{secnumdepth}{0}
\setlength{\parindent}{0pt}
\setlength{\parskip}{5pt}
\setlength{\emergencystretch}{3em}
\begin{document}
\section*{The circular embedding conjecture}\label{3162}
{\small UnsolvedMath ID 3162 · OPG-138 · Graph Theory · OpenGarden}
\subsection{Definitions and mathematical statement}
Let $G$ be a finite simple $2$-connected graph, meaning that
it has at least three vertices and remains connected after
deleting any one vertex. A surface here means a compact
connected topological $2$-manifold without boundary. It may be
orientable or nonorientable, and its genus is not fixed in
advance.

An embedding of $G$ in a surface places vertices at distinct
points and realizes edges as arcs whose interiors are mutually
disjoint and avoid the vertices. It is cellular if each
component of the surface minus the embedded graph is an open
disc; these components are its faces. A facial boundary is
usually described by a closed walk, which may in general
repeat vertices or edges.

The circular embedding conjecture asks whether every such $G$
has a cellular embedding in some surface in which the boundary
walk of every face is a simple cycle. Thus each facial boundary
visits each of its vertices once before returning to its start;
the face closes to a disc bounded by an embedded circle.

The surface and rotation of incident edges may be chosen to
suit $G$. No planar embedding or orientable embedding is
required. Each edge has two sides in a surface, so facial
cycles in such an embedding cover every edge twice when faces
are counted with multiplicity. However, the target includes
the surface condition at vertices: gluing discs along an
arbitrary cycle cover need not give manifold neighbourhoods.
\subsection{Short English statement}
Can every 2-connected graph be embedded on some closed surface so each face is a disc bounded by a simple cycle?
\subsection{Research attempt: a local criterion for gluing cycle discs}
\paragraph{Scope and current literature.}
GPT-6 Astra Ultra, 1 October 2026. The surface may be nonorientable;
the requirement is that vertex neighbourhoods are discs and facial
boundaries are simple cycles. A significant current source is Oum's
2026 exposition [S3] of the announced OpenAI proof of the ordinary
cycle double cover theorem. Its introduction distinguishes the stronger
embedding problem, which remains a separate problem in general, and
explains the cubic consequence. The argument here proves that local
connection explicitly and supplies a counterexample to naive gluing in
higher degree. No proof of the general embedding conjecture is claimed.

\paragraph{1. The link criterion.}
Suppose a connected graph has a list of simple circuits covering each
edge exactly twice. Attach a closed disc along each listed circuit,
identifying its boundary edges and vertices with the corresponding
edges and vertices of the graph. Interior points of a graph edge have
two incident half-disc neighbourhoods, so they have manifold
neighbourhoods. The possible obstruction is at a vertex $v$.

Construct a multigraph $L_v$ with one vertex for each edge incident
with $v$. For each face corner at $v$, join the two incident edges
used by that corner. Every vertex of $L_v$ has degree two, because
each original edge occurs twice in the cover. Moreover there are no
loops: a simple circuit uses two distinct incident edges at a vertex.
The neighbourhood of $v$ is a cone on this link. It is a disc exactly
when $L_v$ is a single circle, allowing a two-edge circle in the
multigraph convention. If every link is connected, the finite quotient
is a closed connected surface and the attached disc interiors are its
faces. Their boundaries remain the prescribed simple circuits.

\paragraph{2. Why cubic vertices cause no obstruction.}
At a cubic vertex, label the incident edges $1,2,3$ and let $a_{ij}$
count corners using edges $i,j$. The double-cover conditions give
\[
 a_{12}+a_{13}=2,\qquad a_{12}+a_{23}=2,\qquad
 a_{13}+a_{23}=2.
\]
Consequently all three corner counts equal one. The link is a triangle
and therefore connected. A circuit double cover of a connected cubic
graph automatically produces the required surface. In particular,
applying the external ordinary double-cover theorem gives the cubic
case; that implication is not an independent proof of the theorem.

There is also an entirely explicit subfamily. In a properly
three-edge-coloured cubic graph, take the three unions of two colour
classes and decompose each into its circuit components. Each edge
belongs to two of the unions. The preceding link calculation gives a
circular embedding directly. Subdividing any edges preserves such an
embedding, because the new degree-two vertices merely subdivide facial
boundary arcs and still have circular links.

\paragraph{3. A concrete failure of unrestricted gluing.}
On vertices $0,1,2,3,4$, the cycles
\[
 C=(0,1,2,3,4,0),\qquad D=(0,2,4,1,3,0)
\]
partition $E(K_5)$. The list $C,C,D,D$ is therefore a circuit double
cover. At every vertex, however, its link is the disjoint union of
two two-edge circles: one joins the two incident $C$-edges and the
other joins the two incident $D$-edges. Thus the glued object is
pinched at every graph vertex. Even four-connectivity of the original
graph does not make this particular cover a surface embedding.
This is an obstruction to the proposed construction, not a claim that
$K_5$ lacks a circular embedding.

\paragraph{Verification and first remaining gap.}
The script [S4] verifies the two edge-disjoint Hamilton cycles of $K_5$
and the two components in each link. The distinction between a list
of circuits and a list of even subgraphs is harmless for ordinary
coverage after circuit decomposition, but connected vertex links must
be checked afterwards. For vertices of degree at least four, the
local parity equations permit several link components. The missing
step is selecting or modifying a global circuit cover so all those
links become connected simultaneously while each face stays simple.
The ordinary double-cover theorem alone does not supply that step.

\subsection{Final status}
UNRESOLVED. A local surface certificate and the cubic reduction are
proved, and a higher-degree gluing obstruction is verified. The cubic
existence consequence is credited to the current external theorem;
this attempt does not resolve arbitrary two-connected graphs.
Completion estimate: no defensible global percentage.

\subsection{Sources}
\begin{itemize}
\item[S1] ULAM.AI and the UnsolvedMath Contributors, supplied problems.json, record 3162 (OPG-138), OpenGarden collection. Catalogue identity and source transcription; CC BY 4.0.
\url{https://huggingface.co/datasets/ulamai/UnsolvedMath}.
\item[S2] Open Problem Garden, The circular embedding conjecture. Original problem and discussion; the mathematical formulation below is expanded from the supplied source record.
\url{http://www.openproblemgarden.org/op/the_circular_embedding_conjecture}.
\item[S3] S.-i. Oum, \emph{A proof of the cycle double cover conjecture by OpenAI: An exposition}, arXiv:2607.16356, version 3 (2026), especially the introduction and discussion of strong embeddings. Checked 1 October 2026.
\url{https://arxiv.org/abs/2607.16356}.
\item[S4] GPT-6 Astra Ultra, exact link and incidence checks, 1 October 2026: \texttt{scripts/check\_attempt\_batch02\_connectivity\_2026\_10\_01.py} in this repository; run with Python 3.
\end{itemize}
\end{document}
